% Elaborado por Fabian Gomez
% Version 1.0

\section{Appendices}\label{sec:appendices}

Esta sección agrupa anexos de apoyo para el SyRS/SRS del Rover Olympus. Incluye (i) la matriz de trazabilidad maestra, (ii) el registro de TBDs y plan de resolución, (iii) el registro de requisitos diferidos/superseded, (iv) la FMEA, (v) los presupuestos de ingeniería (potencia, masa, RF), (vi) el glosario consolidado y (vii) el historial de revisiones. Los anexos se mantienen concisos y referencian, sin duplicar, las definiciones maestras de las secciones de requisitos, V\&V, trazabilidad y arquitectura.

\subsection{Traceability Matrix (SyRS $\rightarrow$ Subsystem $\rightarrow$ V\&V)}
La siguiente tabla resume la trazabilidad de los requerimientos de sistema hacia sus derivados, método de verificación primaria, nivel V\&V y estado de cierre. \textbf{Nivel:} S=Sistema, I=Integración, U=Unitario. \textbf{Estado:} OPEN=pendiente de evidencia; CLOSED=evidencia registrada. \cite{nasa_cm, segoldmine_ppi}

\begin{landscape}
\centering
\scriptsize
\setlength{\tabcolsep}{3pt}
\renewcommand{\arraystretch}{1.3}
\begin{longtable}{| l | p{4.8cm} | p{3.6cm} | l | c | c |}
\caption{Matriz de Trazabilidad Maestra del Rover Olympus (NASA SWE-052).}\label{tab:rtm}\\
\hline
\textbf{ID SyRS} & \textbf{Descripción Resumida} & \textbf{Derivado / Referencia} & \textbf{Método V\&V} & \textbf{Nivel} & \textbf{Estado} \\ \hline
\endfirsthead
\hline
\textbf{ID SyRS} & \textbf{Descripción Resumida} & \textbf{Derivado / Referencia} & \textbf{Método V\&V} & \textbf{Nivel} & \textbf{Estado} \\ \hline
\endhead
\hline
\endfoot
\hline
\endlastfoot

\multicolumn{6}{|l|}{\textbf{Requisitos Funcionales de Sistema (RF)}} \\ \hline
RF-001 & Percepción Visual CSI (FOV según lente OV5647; $\ge$120° con M12 wide-angle) & GNC-REQ-001 & Demonstration/Test & S & OPEN \\
RF-002 & Clasificación de Terreno ($\ge$95\%) & GNC-REQ-002 & Analysis & S & OPEN \\
RF-003 & Evitación Reactiva (HC-SR04 $\le$200\,mm, VL53L0X $\le$150\,mm) & TRAC-REQ-001; GNC-REQ-002 & Test & S & OPEN \\
RF-004 & Compensación de Slip \textit{[SUPERSEDED]} & GNC-REQ-003 & -- & -- & SUPERSEDED \\
RF-004-R1 & Detección Slip (Encoder, error $<$5\% vs ground truth) & GNC-REQ-003; CR-002 & Analysis/Test & S & OPEN \\
RF-005 & Fault Stop ($\le$500\,ms ante stall/proximidad/OC\_FAULT) & PROP-REQ-002 & Demonstration/Test & S & OPEN \\
RF-006 & Integridad TM/CMD (CRC-32 CSP) & COMM-REQ-001 & Analysis/Inspection & S & OPEN \\ \hline

\multicolumn{6}{|l|}{\textbf{Requisitos UC Secundarios (SYS-FUN)}} \\ \hline
SYS-FUN-020 & Registro dinámico 5 waypoints seguros & SW-logic MSM/HLC & Demonstration & S & OPEN \\
SYS-FUN-021 & Retreat ante link lost$>$10\,s o slip persistente HLC (\texttt{SLIP\_STALL\_FRAMES = 2}) & SW-logic MSM/HLC & Demonstration & S & OPEN \\
SYS-FUN-030 & Manual override inmediato & SW-logic MSM & Demonstration & S & OPEN \\
SYS-FUN-031 & Transición autónomo$\rightarrow$manual $\le$500\,ms & CDH-REQ-001 & Test & S & OPEN \\
SYS-FUN-040 & Safe Mode entry \textit{[SUPERSEDED — 040a+040b]} & -- & -- & -- & SUPERSEDED \\
SYS-FUN-040a & Safe Mode por voltaje cualquier banco $<$12\,V (INA3221) & EPS-REQ-001 & Test & S & OPEN \\
SYS-FUN-040b & Safe Mode por latencia HLC $>$5\,s & CDH-REQ-001 & Test & S & OPEN \\
SYS-FUN-041 & Safe Mode: TM salud + motores desenergizados & SYS-FUN-040a/040b & Demonstration & S & OPEN \\
SYS-FUN-050 & Sync y cierre de logs antes de apagado & INF-REQ-001 & Inspection & S & OPEN \\
SYS-FUN-051 & Actuadores en parking antes de deactivación & PROP-REQ-002 & Demonstration & S & OPEN \\ \hline

\multicolumn{6}{|l|}{\textbf{Modos Operacionales (MODE)}} \\ \hline
MODE-001..007 & Modos STB/\allowbreak EXP/\allowbreak AVD/\allowbreak CLB/\allowbreak RET/\allowbreak FLT/\allowbreak SAFE (7 estados \texttt{RoverState} + \texttt{SafetyState} paralelo) & Tab.~\ref{tab:op_modes} & Demonstration & S & OPEN \\ \hline

\multicolumn{6}{|l|}{\textbf{Requisitos No Funcionales y de Desempeño (RNF)}} \\ \hline
RNF-001 & Latencia ciclo $\le$2\,s & CDH-REQ-001 & Test & S & OPEN \\
RNF-002 & Resiliencia recuperación $\ge$90\% & ConOps/SW-logic & Demonstration & S & OPEN \\
RNF-003 & Precisión navegación $<$5\% dist. & GNC-REQ-003; RF-004-R1 & Test/Analysis & S & OPEN \\
RNF-007 & Duración Operacional Mínima $\ge$2\,h & CDH-REQ-002 & Demonstration & S & OPEN \\ \hline

\multicolumn{6}{|l|}{\textbf{Interfaces (ICD)}} \\ \hline
ICD-LLC-001 & Interfaz C\&DH$\leftrightarrow$GNC (UART LLC 115200\,bps) & §7 Tab.~\ref{tab:icd_registry} & Inspection & I & OPEN \\
ICD-COMMS-001 & Interfaz C\&DH$\leftrightarrow$COMMS (CSP/Wi-Fi/UHF) & §7 Tab.~\ref{tab:icd_registry} & Inspection & I & OPEN \\ \hline

\multicolumn{6}{|l|}{\textbf{Requisitos de Calidad y Soporte}} \\ \hline
USA-REQ-002 & Despliegue operativo $\le$15\,min & PKG-REQ-001 & Demonstration & S & OPEN \\
SEC-REQ-001 & Autenticación de comandos & COMM-REQ-001 & Analysis & S & OPEN \\
INF-REQ-001 & Persistencia de logs en NVM & CDH-REQ-002 & Inspection & S & OPEN \\
PHY-REQ-001 & Dimensiones 480$\times$340$\times$320\,mm ($\pm$10\%) & STR-REQ-001 & Inspection & S & OPEN \\

\end{longtable}
\end{landscape}

\subsection{TBD Register (Plan de Resolución)}
Conforme a ISO/IEC/IEEE 29148, este registro identifica los parámetros técnicos aún no cuantificados, asigna un responsable y establece un criterio para su resolución antes de la revisión final de V\&V.

\begin{table}[H]
\centering
\scriptsize
\begin{tabularx}{\textwidth}{| l | X | p{3.0cm} | p{2.2cm} | p{3.8cm} |}
\hline
\textbf{TBD-ID} & \textbf{Descripción del Parámetro} & \textbf{Requisito/Elemento} & \textbf{Responsable} & \textbf{Criterio de Resolución} \\ \hline
TBD-OP-01 & Umbrales de severidad multi-nivel para contingencias (\texttt{severity\_\allowbreak allows\_\allowbreak retreat}, \texttt{severity\_\allowbreak requires\_\allowbreak replan}, \texttt{severity\_\allowbreak requires\_\allowbreak retreat}). Impl.~actual HLC: guarda binaria \texttt{dist\_mm < 300\,mm AND state != CLIMB AND not SafeMode} en \texttt{WaypointTracker.should\_retreat} --- sin clasificación graduada. & \S\ref{subsec:uc01_states}; s09a; s06 & Líder GNC & Análisis de simulación en entorno volcánico + integración ELANav. \\ \hline
TBD-OP-02 & Política de reintentos de enlace (\texttt{retry\_allowed}) & \S\ref{subsec:uc01_states}; s09a & Líder COMM & Caracterización de latencia UHF en campo. \\ \hline
TBD-ENV-01 & Nivel de protección antipartículas (IP4X o superior) & STR-REQ-002, SyRS-100 & Líder STR & Pruebas de infiltración con regolito simulado. \\ \hline
TBD-PER-01 & Ground truth físico para validación encoder (distancia real vs. estimada) & SyRS-013, RF-004-R1; CR-002 & Líder GNC & Definir recorrido de referencia con medición física en sitio análogo. \\ \hline
TBD-DATA-01 & Umbral de almacenamiento crítico (\texttt{storage\_near\_full}) & CDH-REQ-002, §11.1 & Líder C\&DH & Capacidad efectiva de microSD vs bitrate TM. \\ \hline
TBD-LOG-01 & Formato y versión de datasets de misión & §9.5.14 (INF-REQ) & Líder C\&DH & Estandarización conforme a ELANav API. \\ \hline
\end{tabularx}
\caption{Registro maestro de parámetros pendientes (TBD) y plan de resolución.}
\end{table}

\subsection{Deferred and Deleted Requirements Register}\label{subsec:deferred-deleted}

Este registro documenta los requisitos postergados o descartados durante el proceso de diseño, preservando sus IDs para mantener la trazabilidad del baseline. Conforme a NASA-NPR-7123.1C e ISO/IEC/IEEE 29148, los IDs de requisito son inmutables y no se reutilizan. \cite{nasa_cm, segoldmine_ppi}

\textbf{Estados de requisito utilizados en este registro:}
\begin{itemize}
    \item \textbf{[DEFERRED]} — Requisito válido pero postergado a una fase futura; el hardware o contexto operacional no está disponible en el scope actual del TFG.
    \item \textbf{[DELETED]} — Requisito descartado del scope; ID congelado para trazabilidad histórica.
    \item \textbf{[EXTERNAL]} — Requisito cuya verificación depende de hardware físico externo al software de vuelo (LLC/HLC). Marcados con \texttt{[EXTERNO]} en §9; estado activo, no forman parte de este registro.
\end{itemize}

\begin{table}[H]
\centering
\scriptsize
\begin{tabularx}{\textwidth}{| l | l | X | l |}
\hline
\textbf{ID} & \textbf{Estado} & \textbf{Razón} & \textbf{Ref. CR} \\ \hline

RF-004 & SUPERSEDED por RF-004-R1 &
Requisito original especificaba compensación de slip mediante EKF. EKF descartado del scope del TFG: implementación de fusión sensorial fuera del alcance actual y hardware IMU no incluido en la configuración de vuelo. Reemplazado por RF-004-R1 (odometría encoder, 6 ruedas, error $<$5\% vs. ground truth). &
CR-002 \\ \hline

COMM-REQ-003 & DEFERRED — FASE-2 &
Hardware RF disponible (Kenwood TM-D710GA + NinoTNC 9600A) pero integración y campaña V\&V postergadas fuera del scope del TFG actual. Candidato para campaña de validación extendida. &
CR-001 \\ \hline

COMM-REQ-004 & DEFERRED — FASE-2 &
Ídem COMM-REQ-003. La integridad multicapa en UHF requiere el enlace físico activo; no verificable en campaña actual. &
CR-001 \\ \hline

COMM-REQ-005 & DEFERRED — FASE-2 &
Ídem COMM-REQ-003. La conmutación automática de ruta requiere al menos dos nodos RF activos (rover + relay) no disponibles en esta fase. &
CR-001 \\ \hline

\end{tabularx}
\caption{Registro de requisitos diferidos y eliminados del Rover Olympus.}
\label{tab:deferred}
\end{table}

\textbf{Nota sobre requisitos [EXTERNO]:} Los requisitos marcados \texttt{[EXTERNO]} en \S\ref{sec:requirements} (STR-REQ-001, TRAC-REQ-002, CDH-REQ-002, EPS-REQ-002b, COMM-REQ-002, GNC-REQ-003) son requisitos \textbf{activos} cuya verificación depende del hardware físico o de documentación de proveedor, no del software de vuelo. Su estado es ACTIVE/EXTERNAL y no forman parte de este registro.

\emph{Nota sobre consolidaciones recientes:} STR-REQ-002 ya no se considera externo --- consolidado con ENV-REQ-003 que aporta método de test (blow test IP4X); EPS-REQ-002 se dividió en EPS-REQ-002a (protección sw del LLC, no externo) y EPS-REQ-002b (protección hw por celda, externo).

\subsection{FMEA — Failure Mode and Effects Analysis}\label{subsec:fmea}

Este análisis identifica los modos de falla relevantes de los componentes y funciones críticas del Rover Olympus, sus efectos sobre la misión UC-01, y las mitigaciones ya implementadas en el software de vuelo (LLC/HLC). El objetivo es demostrar que el diseño considera el comportamiento ante falla (\textit{fail-safe}) y que los requisitos RF-005, RNF-002 y SYS-FUN-040 están respaldados por mecanismos concretos. Análisis conforme a ECSS-Q-ST-30C \cite{ecss_q_st_30}.

\textbf{Escalas conforme a ECSS-Q-ST-30C} \cite{ecss_q_st_30}\textbf{:}
\begin{itemize}
    \item \textbf{Severidad:} \textit{Critical} = pérdida de misión o daño al sistema; \textit{Major} = degradación significativa, misión comprometida; \textit{Minor} = función degradada, misión continúa; \textit{Negligible} = sin impacto operacional.
    \item \textbf{Ocurrencia:} \textit{Frecuente} ($>$1/h); \textit{Probable} (1/h–1/día); \textit{Ocasional} (1/día–1/campaña); \textit{Remota} (1/campaña–1/año); \textit{Improbable} ($<$1/año).
    \item \textbf{Detección:} mecanismo por el cual la falla se detecta antes de causar el efecto indicado.
\end{itemize}

\begin{landscape}
\centering
\scriptsize
\setlength{\tabcolsep}{3pt}
\renewcommand{\arraystretch}{1.3}
\begin{longtable}{| l | p{2.4cm} | p{2.8cm} | c | c | p{2.6cm} | p{3.8cm} | l |}
\caption{FMEA del Rover Olympus — modos de falla, severidad, ocurrencia y mitigaciones (ECSS-Q-ST-30-02C).}
\label{tab:fmea}\\
\hline
\textbf{Componente} & \textbf{Modo de Falla} & \textbf{Efecto en Misión} & \textbf{Clase} & \textbf{Ocurr.} & \textbf{Detección} & \textbf{Mitigación Implementada} & \textbf{Req.} \\ \hline
\endfirsthead
\hline
\textbf{Componente} & \textbf{Modo de Falla} & \textbf{Efecto en Misión} & \textbf{Clase} & \textbf{Ocurr.} & \textbf{Detección} & \textbf{Mitigación Implementada} & \textbf{Req.} \\ \hline
\endhead
\hline
\endfoot
\hline
\endlastfoot

Cámara CSI \newline (VisionSource) &
Timeout sin frame ($>$2\,s) &
Percepción ciega; rover no puede clasificar terreno &
Critical & Ocasional &
Timeout counter en VisionSource &
Timeout $\rightarrow$ MSM transición a AVOIDANCE / FAULT &
RF-001 \newline RF-005 \\ \hline

Encoders ($\times$6) &
Señal perdida o pulsos fuera de rango &
Odometría inválida; slip no detectable &
Major & Remota &
Stall counter LLC por rueda &
Stall detection en 6 encoders $\rightarrow$ SafeMode &
RF-004-R1 \newline RF-005 \\ \hline

Sensor HC-SR04 &
Sin eco / lectura saturada &
Obstáculo no detectado; riesgo colisión &
Critical & Ocasional &
Valor fuera de rango {[}0--400{]}\,cm &
Lectura inválida $\rightarrow$ obstáculo presente (\textit{fail-safe}); VL53L0X y TF02 como respaldo &
RF-003 \\ \hline

TF02 LiDAR \newline (largo alcance) &
Frame USART2 corrupto / FIFO overflow / SIG bajo ($<$7) &
Distancia frontal de largo alcance no fiable; objetos lejanos no anticipados &
Major & Ocasional &
ISR USART2 con ring buffer; validación SIG $\ge$ 7 en HLC &
Frame inválido descartado; LLC continúa con HC-SR04 + VL53L0X como respaldo &
RF-001 \newline RF-003 \\ \hline

VL53L0X \newline (ToF corto alcance) &
Fallo Soft-I$^2$C / lectura saturada ($>$2000\,mm) &
Detección corto alcance no fiable; obstáculo cercano no detectado &
Critical & Remota &
Soft-I$^2$C timeout / rango fuera de [10, 2000]\,mm &
Lectura inválida $\rightarrow$ obstáculo presente (\textit{fail-safe}); HC-SR04 como respaldo &
RF-003 \\ \hline

INA3221 \newline (monitor multi-banco EPS) &
Fallo Soft-I$^2$C / lectura inválida en uno o más bancos &
Batería sin monitoreo; riesgo sobrecarga/sobre-descarga &
Major & Remota &
Lectura Soft-I$^2$C inválida o fuera de rango físico (0--26\,V por banco) &
Valor inválido $\rightarrow$ alerta HLC; voltaje de cualquier banco $<$12\,V $\rightarrow$ SafeMode (SYS-FUN-040a) &
EPS-REQ-001 \newline SYS-FUN-040a \\ \hline

BTS7960 \newline (drivers motores CR/CL/RR/RL) &
Sobrecorriente ($>I_{stall}$ NFP-5840 $\approx$ 6.5\,A) o sobre-temperatura del driver &
Motor bloqueado o driver destruido; pérdida de tracción &
Critical & Ocasional &
ACS712 ADC + umbral \texttt{OC\_FAULT} en LLC (\texttt{config.rs}) &
\texttt{OC\_FAULT} $\rightarrow$ desenergización inmediata del canal afectado + transición a FLT &
RF-005 \newline EPS-REQ-001 \\ \hline

LM335 / NTC \newline (sensores térmicos) &
Lectura errónea / desconexión &
Temperatura no monitoreada &
Minor & Remota &
Valor fuera de rango físico esperado &
ThermalMonitor registra anomalía en log; umbral $>$45\,°C $\rightarrow$ alerta &
RNF-004 \\ \hline

Enlace CSP \newline (TM/CMD) &
Paquete con CRC inválido &
Comando descartado; posible pérdida de telemetría &
Critical & Ocasional &
CRC-32 mismatch en \texttt{unpack()} &
CRC-32 drop; no-ACK $\rightarrow$ retransmisión &
RF-006 \newline COMM-REQ-001 \\ \hline

OBC RPi5 \newline (ciclo HLC) &
Ciclo $>$2\,s (\texttt{cycle\_warn\_ms}) &
Latencia excedida; respuesta tardía a contingencias &
Major & Ocasional &
Timer \texttt{cycle\_warn\_ms} excedido &
\texttt{cycle\_warn\_ms=1500} watchdog; log + alerta GCS &
RNF-001 \newline CDH-REQ-001 \\ \hline

LLC Arduino Mega &
UART timeout / desconexión &
Sin actuación de motores; rover inmovilizado &
Critical & Remota &
Heartbeat ausente $>N$ ciclos &
Heartbeat LLC$\rightarrow$HLC; ausencia $\rightarrow$ FAULT en MSM &
RF-005 \newline SYS-FUN-040b \\ \hline

MSM \newline (state machine) &
Transición de estado inválida &
Comportamiento impredecible del sistema &
Critical & Improbable &
Transición rechazada por guard en Rust &
Validación en \texttt{state\_machine/mod.rs}; \textit{default safe} $\rightarrow$ STB &
RF-005 \newline SyRS-060 \\ \hline

\end{longtable}
\end{landscape}

\textbf{Nota:} Todos los modos de falla Critical/Major convergen en SafeMode (SYS-FUN-040a/040b) o FAULT en la MSM. No se identificaron modos de falla sin mitigación en los componentes de vuelo actuales. Clasificación de severidad y ocurrencia conforme a ECSS-Q-ST-30C \cite{ecss_q_st_30}.

% ============================================================
\subsection{Engineering Budgets}\label{subsec:eng_budgets}
% ============================================================

Este apéndice establece los presupuestos de ingeniería (potencia, masa y enlace RF) del Rover Olympus con valores derivados de hojas de datos del fabricante (\textbf{[DS]}), estimaciones de ingeniería (\textbf{[EST]}) y parámetros pendientes de medición en hardware real (\textbf{[TBM]}). Los valores [TBM] deben resolverse antes del cierre de la campaña V\&V para cerrar los requisitos RNF-007 (MOD $\ge$2\,h) y EPS-REQ-001 (monitoreo multi-banco INA3221) \cite{rpi5_brief, arduino_mega_ds}.

% ------------------------------------------------------------
\subsubsection{Power Budget (Bus 12 V — 3 bancos 18650 NMC)}
% ------------------------------------------------------------

\begin{landscape}
\centering
\scriptsize
\setlength{\tabcolsep}{3pt}
\renewcommand{\arraystretch}{1.3}
\begin{longtable}{| p{2.0cm} | p{4.5cm} | c | r | r | r | r | r | c |}
\caption{Presupuesto de potencia del Rover Olympus (modo traversal, bus 12\,V, 3 bancos 18650 NMC).
  Valores [TBM] requieren banco de motores con corriente medida por ACS712 por canal (INA3221 mide voltaje multi-banco). Criterio RNF-007: energía total $\ge$89.4\,Wh
  $\Rightarrow$ capacidad combinada $\ge$5\,400\,mAh ($\ge$80\,Wh con 10\,\% derating).}
\label{tab:power_budget}\\
\hline
\textbf{Subsistema} & \textbf{Componente} & \textbf{V\textsubscript{bus}} &
\textbf{I\textsubscript{idle} (A)} & \textbf{P\textsubscript{idle} (W)} &
\textbf{I\textsubscript{peak} (A)} & \textbf{P\textsubscript{peak} (W)} &
\textbf{E\textsubscript{2h} (Wh)} & \textbf{Fuente} \\ \hline
\endfirsthead
\hline
\textbf{Subsistema} & \textbf{Componente} & \textbf{V\textsubscript{bus}} &
\textbf{I\textsubscript{idle} (A)} & \textbf{P\textsubscript{idle} (W)} &
\textbf{I\textsubscript{peak} (A)} & \textbf{P\textsubscript{peak} (W)} &
\textbf{E\textsubscript{2h} (Wh)} & \textbf{Fuente} \\ \hline
\endhead
\hline
\endfoot
\hline
\endlastfoot

C\&DH  & RPi5 (OBC/HLC)         & 5\,V  & 0.636 & 3.18  & 1.506 & 7.53  & 6.36  & [DS] \cite{rpi5_brief} \\ \hline
C\&DH  & Arduino Mega (LLC)      & 5\,V  & 0.050 & 0.25  & 0.050 & 0.25  & 0.50  & [DS] \cite{arduino_mega_ds} \\ \hline
GNC    & Cámara OV5647 (CAM1)    & 3.3\,V & 0.150 & 0.50  & 0.150 & 0.50  & 1.00  & [DS] \\ \hline
GNC    & HC-SR04 ($\times$1)     & 5\,V  & 0.003 & 0.015 & 0.003 & 0.015 & 0.03  & [DS] \\ \hline
EPS    & INA3221 (monitor multi-banco) & 3.3\,V & 0.001 & 0.003 & 0.001 & 0.003 & 0.01 & [DS] \\ \hline
EPS    & ACS712 ($\times$6, sense corriente motor) & 5\,V & 0.060 & 0.30 & 0.060 & 0.30 & 0.60 & [DS] \\ \hline
PROP   & 6$\times$ Motor NFP-5840 & 12\,V & {[TBM]} & {[TBM]} & {[TBM]} & $\approx$14.4 & {[TBM]} & [EST] \\ \hline
COMM   & WiFi 802.11n (RPi5 integ.) & incl. & -- & -- & -- & $\le$1.0 & $\le$2.0 & [EST] \\ \hline
\hline
\multicolumn{2}{|l|}{\textbf{Total estimado — modo traversal (nominal)}}
  & --    & --    & $\approx$18   & --    & $\approx$44.7 & $\approx$89.4 & [EST] \\ \hline

\end{longtable}
\end{landscape}

% ------------------------------------------------------------
\subsubsection{Mass Budget}
% ------------------------------------------------------------

\begin{table}[H]
\centering
\scriptsize
\renewcommand{\arraystretch}{1.3}
\begin{tabularx}{\textwidth}{| l | l | r | c | X |}
\hline
\textbf{Subsistema} & \textbf{Componente} & \textbf{Masa (g)} & \textbf{Fuente} & \textbf{Notas} \\ \hline

C\&DH  & RPi5 (SBC)                & 45    & [DS]  & PCB sin carcasa \\ \hline
C\&DH  & Arduino Mega 2560         & 37    & [DS]  & PCB sin carcasa \\ \hline
C\&DH  & microSD 64\,GB (clase 10) & 2     & [DS]  & \\ \hline
GNC    & Cámara OV5647 (CAM1)      & 6     & [DS]  & módulo + lente seleccionado (ver GNC-REQ-001) \\ \hline
GNC    & HC-SR04                   & 8     & [DS]  & \\ \hline
EPS    & INA3221 breakout (triple) & 3     & [DS]  & monitoreo multi-banco \\ \hline
EPS    & ACS712 ($\times$6, current sense) & 30 & [DS] & 5\,g/módulo, uno por motor \\ \hline
EPS    & LM335 + NTC ($\times$2)   & 5     & [EST] & sensores térmicos \\ \hline
PROP   & 6$\times$ Motor NFP-5840  & 2\,160 & [EST] & $\approx$360\,g/motor (estimado) \\ \hline
PROP   & Controladores de motor    & 120   & [EST] & 6$\times$ driver \\ \hline
COMM   & Módulo WiFi integrado     & --    & [DS]  & incluido en RPi5 \\ \hline
STR    & Chasis / estructura       & {[TBM]} & [TBM] & Medición directa con balanza \\ \hline
STR    & Ruedas ($\times$6)        & {[TBM]} & [TBM] & \\ \hline
STR    & Arnés / cableado          & {[TBM]} & [TBM] & Estimado $\approx$200\,g \\ \hline
EPS    & 18650 NMC ($\times$3 bancos) & {[TBM]} & [TBM] & celdas + portabaterías + protección \\ \hline
\hline
\multicolumn{2}{|l|}{\textbf{Electrónica conocida (excl. estructura)}} & $\approx$2\,386 & [DS/EST] & sin arnés ni batería \\ \hline
\multicolumn{2}{|l|}{\textbf{Total estimado (sistema completo)}}       & 5\,000--7\,000 & [EST] & pendiente medición estructura \\ \hline
\end{tabularx}
\caption{Presupuesto de masa del Rover Olympus. Componentes [TBM] deben medirse con balanza de precisión antes de campaña.}
\label{tab:mass_budget}
\end{table}

% ------------------------------------------------------------
\subsubsection{RF Link Budget (2.4\,GHz, 50\,m LoS)}
% ------------------------------------------------------------

La siguiente tabla aplica la ecuación de Friis \cite{friis_1946} para verificar el cierre del enlace WiFi 802.11n \cite{ieee_802_11n} en la distancia operacional máxima declarada en COMM-REQ-001 (50\,m, LoS). \cite{omg_sysml_mbse, segoldmine_ppi}

\begin{table}[H]
\centering
\scriptsize
\renewcommand{\arraystretch}{1.3}
\begin{tabularx}{\textwidth}{| l | r | c | X |}
\hline
\textbf{Parámetro} & \textbf{Valor} & \textbf{Unidad} & \textbf{Fuente / Notas} \\ \hline

Frecuencia $f$                         & 2\,400  & MHz  & Canal 1--11 WiFi 802.11n \\ \hline
Distancia $d$                          & 50      & m    & COMM-REQ-001 (máx. LoS) \\ \hline
Potencia transmisión $P_t$             & +20     & dBm  & Límite legal WiFi ISM (típico 100\,mW) \\ \hline
Ganancia antena TX $G_t$               & +2      & dBi  & Antena dipolo integrada (RPi5) \\ \hline
Ganancia antena RX $G_r$               & +2      & dBi  & Antena dipolo integrada (GCS laptop) \\ \hline
FSPL a 50\,m, 2.4\,GHz                & 74.2    & dB   & $\text{FSPL}=20\log_{10}(d)+20\log_{10}(f)$\newline${}+20\log_{10}(\tfrac{4\pi}{c})$ \cite{friis_1946} \\ \hline
Pérdidas adicionales $L_{misc}$        & +5      & dB   & Margen de desvanecimiento, conectores \\ \hline
\hline
\textbf{Potencia recibida} $P_r$       & \textbf{$-$55.2}  & dBm  & $P_r = P_t + G_t + G_r - \text{FSPL} - L_{misc}$ \\ \hline
Sensibilidad del receptor              & $-91$   & dBm  & IEEE 802.11n MCS0 (6\,Mbps) \cite{ieee_802_11n} \\ \hline
\textbf{Margen de enlace}              & \textbf{$+$35.8} & dB   & $M = P_r - S_{min}$; criterio aceptable $\ge$10\,dB \\ \hline
\textbf{Estado}                        & \multicolumn{2}{c|}{\textbf{CERRADO}} & Margen holgado; COMM-REQ-001 satisfecho en LoS \\ \hline
\end{tabularx}
\caption{Presupuesto de enlace RF WiFi 2.4\,GHz a 50\,m (ecuación de Friis). El margen de $+$35.8\,dB permite
  absorber obstáculos parciales sin degradar el enlace por debajo de la sensibilidad del receptor. El enlace UHF
  (COMM-REQ-002/003) está diferido a Fase-2 (CR-001).}
\label{tab:rf_link_budget}
\end{table}

\textbf{Nota sobre [TBM]:} Los valores de corriente de motores (banco de pruebas), masa de estructura/ruedas y RSSI Wi-Fi en sitio análogo deben medirse durante la campaña de integración y hardware (§11 V\&V). Estos cierres completan los requisitos EPS-REQ-001, STR-REQ-001 y COMM-REQ-001 respectivamente.

\subsection{Glossary}

Este glosario consolida términos clave utilizados en el SyRS/SRS. Para definiciones operativas completas, ver \S\ref{sec:terms} (Terms, Definitions, Acronyms), que es la fuente autoritativa; este apéndice solo recoge un subconjunto frecuente para referencia rápida.

\textbf{Términos clave (consolidación):}
\begin{itemize}
    \item \textbf{ELANav:} Biblioteca embebida para la navegación de sistemas espaciales de misión crítica. Proyecto de investigación TEC (Escuela de Ingeniería Electrónica, coord.\ Johan Carvajal Godínez, 2025--2027). Olympus es la plataforma rover de validación de esta biblioteca.
    \item \textbf{Rover Olympus:} Rover académico tipo exploración, plataforma de validación de ELANav (7 subsistemas).
    \item \textbf{UC-01:} Exploración de superficie (caso núcleo); navegación autónoma continua con percepción y evitación.
    \item \textbf{Sitio análogo:} Terreno volcánico / regolito simulado usado para validación (pendientes hasta 20°, obstáculos, transiciones de suelo).
    \item \textbf{Slip ratio / deslizamiento:} Diferencia entre velocidad teórica del eje motriz y velocidad real de avance del rover; en Olympus la detección es \emph{event-based} (stall de encoder) y no porcentual, ver RF-004-R1.
    \item \textbf{GCS / Estación terrena:} Sistema externo operado por Rover Planner para uplink/downlink.
    \item \textbf{CSP:} CubeSat Space Protocol (encapsulamiento TM/CMD).
    \item \textbf{V\&V:} Verification \& Validation, con métodos T/A/D/I (Test/Analysis/Demonstration/Inspection).
    \item \textbf{BDD/IBD/PBS/N2/ASM:} Artefactos de arquitectura (SysML y descomposición física) usados para trazabilidad y soporte de integración.
\end{itemize}

\subsection{Revision History}

El historial de cambios del SyRS/SRS Rover Olympus se mantiene en el \emph{Change Register} (\S\ref{subsec:change-register}, Tab.~\ref{tab:change-register}), que documenta los Change Requests (CR-001, CR-002, CR-003\dots) aprobados sobre el baseline, sus motivos, los Configuration Items afectados y la decisión del mini-CCB. Cada CR se enlaza, donde aplica, a los commits en el repositorio del proyecto.

\emph{Convención:} los IDs de requisito son inmutables conforme a NASA-NPR-7123.1C; cambios sustanciales se gestionan vía \emph{supersession} (p.\,ej., RF-004 $\rightarrow$ RF-004-R1) en lugar de renombrar IDs existentes.
